\documentclass[12pt,letterpaper,oneside]{article}

\input{./latexGoodPractices/preamble}

%==============================================================
% FILL THIS SECTION

\newcommand{\reportTitle}{Random maths}
\newcommand{\reportAuthors}{Nicolas Lauzon, Fran\c{c}ois Pomerleau}
\newcommand{\reportDate}{\today} % or manually: November 23, 2017

\newcommand{\reportVersions}{
  0.1 & October 7, 2026 & First draft %\\
  %1.0 & November 23, 2017 & Final writing
}
% Change to your specific bibliography file
\addbibresource{./latexGoodPractices/exampleReferences.bib}
%\addbibresource{./references.bib}
%==============================================================

% ---------------------------------------------------------------
% Load style
%----------------------------------------
% Page style

% Set the page size
\addtolength{\hoffset}{-0.75in} \addtolength{\voffset}{-0.75in}
\setlength{\textwidth}{7in} \setlength{\textheight}{8.25in}
\setlength{\headheight}{0.6in}
\setlength{\headsep}{0.2in}

\setlength{\footskip}{40pt}
\setlength{\fboxsep}{12pt}

% Set the paragraph skip
\setlength{\parskip}{3pt}

% Access to a counter for the number of pages
\usepackage{lastpage}

% To allow text justify on the right
\usepackage{ragged2e}

%----------------------------------------
% Title style

\newcommand{\makeCustomTitle}
{
\begin{center}
\LARGE{\textbf{\reportTitle{}}}
\\
\vspace{5pt}
\normalsize{\reportAuthors{}}
\\
\reportDate{}
\end{center}
\begin{flushright}
\footnotesize{
\begin{tabu}{ccX}
\toprule
\emph{Version} & \emph{Date} & \emph{Comments} \\
\midrule
\reportVersions{} \\
\bottomrule
\end{tabu}
}
\end{flushright}
}

%----------------------------------------
% Section style
\usepackage{sectsty}

% Set the section labeling font
\allsectionsfont{\textsf\bfseries}

%----------------------------------------
% Caption style
\usepackage[font=small, labelfont=bf, skip=5pt]{caption}

%----------------------------------------
% header style
\usepackage{fancyhdr}



% Define the title page style
\fancypagestyle{titlePage}{%
\fancyhf{}%

\fancyhead[L]{\includegraphics[height=0.45in]{UL_N}}
\fancyhead[C]{\raisebox{0.2in}{\textsc{Technical Report}}}
\fancyhead[R]{\includegraphics[height=0.45in]{norlab_logo_acronym_dark}}
\fancyfoot[C]{\thepage/\pageref*{LastPage}}

\renewcommand{\headrulewidth}{0.1pt}
\renewcommand{\footrulewidth}{0.2pt}
}

% Define the page style for the other pages
\fancypagestyle{plain}{%
\fancyhf{}
\fancyhead[L]{\reportTitle{}}
\fancyfoot[C]{\thepage/\pageref*{LastPage}}
\renewcommand{\headrulewidth}{0.1pt}
\renewcommand{\footrulewidth}{0.1pt}
}

% Set the page style for all the document except the first page
\pagestyle{plain}

%----------------------------------------
% footnote style
\usepackage{fnpos}
% Fix the footnotes location
\makeFNbottom \makeFNbelow

% ---------------------------------------------------------------
% Author

\author{\reportAuthors{} \\
       Laval University\\
       1065, av.\ de la Médecine \\
       Quebec, Qc \\
       Canada G1V 0A6 \\
}

% ---------------------------------------------------------------
% PDF setup
\hypersetup{%
    pdftitle={\reportTitle},
    pdfauthor={\@author},
    pdfkeywords={research, project, robotics, norlab, Northern Robotics Lab, Master's},
    pdfsubject={},
    pdfstartview={},
    urlcolor=cyan,
    linkcolor=red,
}%
% ---------------------------------------------------------------

% Add your other packages here
% make sure to check prembule.tex first!
\usepackage{tikz}
\usetikzlibrary{positioning, arrows.meta, matrix}

% Robot states, each with its own color so it is easy to track across equations
\colorlet{colP}{ForestGreen}
\colorlet{colR}{BrickRed}
\colorlet{colV}{RoyalPurple}
\colorlet{colW}{BurntOrange}
\colorlet{colA}{TealBlue}
% The symbol and its frames are colored; dots, transposes and time subscripts stay black.
% Optional argument is the time subscript, e.g. \Bv[t+\Delta t]
\newcommand{\Gp}[1][]{\textcolor{colP}{{}^{\mathcal{G}}\bm{p}}_{#1}}
\newcommand{\Gpdot}{\textcolor{colP}{{}^{\mathcal{G}}}\dot{\textcolor{colP}{\bm{p}}}}
\newcommand{\BR}[1][]{\textcolor{colR}{{}^{\mathcal{B}}_{\mathcal{G}}\bm{R}}_{#1}}
\newcommand{\BRdot}{\textcolor{colR}{{}^{\mathcal{B}}_{\mathcal{G}}}\dot{\textcolor{colR}{\bm{R}}}}
\newcommand{\BRT}[1][]{\textcolor{colR}{{}^{\mathcal{B}}_{\mathcal{G}}\bm{R}}_{#1}^T}
\newcommand{\Bv}[1][]{\textcolor{colV}{{}^{\mathcal{B}}\bm{v}}_{#1}}
\newcommand{\Bvdot}{\textcolor{colV}{{}^{\mathcal{B}}}\dot{\textcolor{colV}{\bm{v}}}}
\newcommand{\Bw}[1][]{\textcolor{colW}{{}^{\mathcal{B}}\bm{\omega}}_{#1}}
\newcommand{\Ba}[1][]{\textcolor{colA}{{}^{\mathcal{B}}\bm{a}}_{#1}}
\DeclareMathOperator{\Exp}{Exp}
\DeclareMathOperator{\Log}{Log}

%================================================================
\begin{document}
\makeCustomTitle
\thispagestyle{titlePage}

% ---------------------------------------------------------------
\begin{abstract}
  Describe in 3--4 sentences the goal of this technical report.
  \lightlipsum[1]
\end{abstract}

% ---------------------------------------------------------------
\section{Definitions}

Use what you need in the symbols and functions defined in \autoref{tab:symbols}, \autoref{tab:pointCloudSymbols}, and \autoref{tab:functions}.
Stay consistent with the lab's notation and your own symbols.

\begin{table}[htbp]
  \centering
  \caption{General symbol definitions.}
  \label{tab:symbols}
  \begin{tabu}{cX}

    \toprule
    \emph{Symbol} & \emph{Explanation} \\

    \midrule
    $a$ & A scalar \\
    $\bm{v}$ & A vector of $D$ dimensions, $\bm{v} = {[v_1, v_2, \ldots, v_d]}^T$ \\
    $\mathcal{S}$ & A set of vectors with $N$ elements in it, $\mathcal{S} = \{\bm{v}_1, \bm{v}_1, \ldots, \bm{v}_n\}$\\
    $\bm{S}$ & Matrix representation of a set $\mathcal{S}$ of size $D \times N$, $\bm{S} = [\bm{v}_1, \bm{v}_1, \ldots, \bm{v}_n]$\\
    $\bm{a} \cdot \bm{b} = c$ & Dot product of two vectors, also known as inner product \\
    $\bm{a}^T \bm{b} = c$ & Inner product of two vectors, also known as dot product \\
    $\mathrm{f}(\cdot)$ & A function, short version of $\mathrm{function}(\cdot)$, same font as $\log$, $\exp$, $\arg\min$\\
    $\|\bm{v}\|_2$ & Euclidian distance or $\ell^2$-distance, $\|\bm{v}\|_2 = \sqrt{\bm{v}^T \bm{v}}$ \\

    \bottomrule

  \end{tabu}
\end{table}

\begin{table}[htbp]
  \centering
  \caption{Symbol definitions for point clouds and registration.}
  \label{tab:pointCloudSymbols}
  \begin{tabu}{cX}

    \toprule
    \emph{Symbol} & \emph{Explanation} \\
    \midrule

    $d$ & Indices used for the dimension of both the point clouds with $d = \{1,2, \ldots, D\}$ \\
    $i$ & Indices used for a point in the reading (moving) point cloud with $i = \{1,2, \ldots, I\}$\\
    $\bm{p}_i$ & A point in the reading point cloud, $\bm{p} = {[x, y, z]}^T$ for a 3D point\\
    $\mathcal{P}$ & A set of discrete sample points representing the reading (moving) point cloud,  $\mathcal{P} = \{\bm{p}_1, \bm{p}_1, \ldots, \bm{p}_i\}$ \\
    $\bm{P}$ & Matrix version of $\mathcal{P}$ with a size of $D \times I$ \\

    $j, q_j, \mathcal{Q}, \bm{Q}$ & Same definitions but for the reference (static) point cloud \\

    $k$ & Indices used when points are matched (points in $\bm{p}$ matched to $\bm{q}$) with $k = \{1,2, \ldots, K\}$\\
    $\bm{e}_k$ & Matched error between two points, $\bm{e}_k = \bm{p}'_i - \bm{q}_j$ \\
    $\mathcal{E}, \bm{E}$ & Set and its matrix version of size $D \times K$ containing all matched error \\

    $\bm{x}$ & States of the robot \\
    $\bm{x}_k$ & States of the robot at the same time as the point $\bm{p}_k$ \\
    $\bm{p}'_k$ & Moved reading point cloud as  $\bm{p}'_k = \mathrm{T}(\bm{x}, \bm{p})$ \\
    $\bm{R}$ & Rotation matrix of size $D \times D$ \\
    $\bm{T}$ & Rigid transformation matrix of size $(D+1) \times (D+1)$ \\
    $\bm{n}_{q_k}$ & Normal vector representing a plane at $\bm{p}_k$ \\
    $\bm{W}_{\bm{q}_k}$ & Covariance matrix of size $D \times D$ expressed in the reference frame of $\bm{q}_k$\\
    \bottomrule

  \end{tabu}
\end{table}

\begin{table}[htbp]
  \centering
  \caption{Function Definitions.}
  \label{tab:functions}
  \begin{tabu}{cX}

    \toprule
    \emph{Function} & \emph{Explanation} \\
    \midrule

    $\mathrm{T}(\bm{x}, \bm{p})$ & Function transforming moving the reading point cloud \\
    $\mathrm{T}(\bm{x}, \bm{p}_k)$ & Rigid transformation with all points in the reading using the same states\\
    $\mathrm{T}(\bm{x}_k, \bm{p}_k)$ & Flexible transformation with each point moved by their own states\\
    $\mathrm{R}(\bm{x})$ & Function building a rotation matrix from the states \\
    $\mathrm{J}(\bm{x}) = a$ & Objective function also known as error function, loss function, cost function\\
    $\mathrm{match}(\mathcal{P}, \mathcal{Q}) = \mathcal{E}$ & Matching function generating a set of error vectors $\mathcal{E}$ with $K$ elements in it\\

    \bottomrule

  \end{tabu}
\end{table}

\newpage

\section{Maths}

\subsection{Relation between a unit-speed curve and a trajectory in the time domain}
A unit-speed curve $\gamma(s)$ is parametrized by arc length $s$.
Formally, we define a unit-speed curve by:
\begin{align}
  \left\| \frac{d\gamma(s)}{ds} \right\| & = 1, \quad \forall s, \\
  \mathbf{T}(s) & = \frac{d\gamma(s)}{ds}, \\
  \kappa(s) & = \left\| \frac{d\mathbf{T}(s)}{ds} \right\|, \\
  \mathbf{N}(s) & = \frac{1}{\kappa(s)} \frac{d\mathbf{T}(s)}{ds}, \quad \kappa(s) > 0.
\end{align}
Here, $\mathbf{T}$ is the unit tangent vector to the curve.
Since $\mathbf{T}$ has constant length, it can only rotate, so $\frac{d\mathbf{T}}{ds}$ is perpendicular to $\mathbf{T}$ and $\mathbf{N}$ is a unit vector perpendicular to $\mathbf{T}$.
The curvature $\kappa$ is a magnitude, $\mathbf{N}$ is its direction.
Intuitively, curvature is how much a curve deviates from being a straight line [wikipedia].

A trajectory $r(t)$ is the position of the robot as a function of time.
We can rewrite it as $r(t) = \gamma(s(t))$ where $s(t)$ is the arc-length profile.
The speed is the rate of change of the arc length,
\begin{equation}
  v(t) = \frac{ds}{dt}.
\end{equation}
A unit-speed curve is traversed at constant speed, so a change of acceleration can only change the heading.
By adding time here, the speed is no longer constant and we have a new component of the acceleration that is tangent to the curve.
Tangential acceleration modifies the speed and perpendicular (normal) acceleration modifies the heading.
Formally, using the chain rule and the product rule, we have:
\begin{align}
  \frac{d r(t)}{dt} & = \frac{d\gamma(s)}{ds} \frac{ds}{dt} = v\, \mathbf{T}(s(t)), \\
  \frac{d^2 r(t)}{dt^2} & = \frac{dv}{dt}\, \mathbf{T} + v \frac{d\mathbf{T}(s)}{ds} \frac{ds}{dt} = \frac{dv}{dt}\, \mathbf{T} + \kappa v^2 \mathbf{N}.
\end{align}
The term $\frac{dv}{dt}\,\mathbf{T}$ is the tangential acceleration and $\kappa v^2 \mathbf{N}$ is the centripetal acceleration, which vanishes on a straight line ($\kappa = 0$).

\subsection{State relations}

The states of the robot are linked by three kinematic equations,
\begin{align}
  \Gpdot & = \BRT \, \Bv,                    \label{eq:kinPos} \\
  \Bvdot & = \Ba - \Bw \times \Bv,           \label{eq:kinVel} \\
  \BRdot & = -{\Bw}^\wedge \, \BR,           \label{eq:kinRot}
\end{align}
where ${(\cdot)}^\wedge$ builds the skew-symmetric matrix such that $\bm{a}^\wedge \bm{b} = \bm{a} \times \bm{b}$.
Left superscripts give the frame, global $\mathcal{G}$ or body $\mathcal{B}$, and $\BR$ rotates vectors from $\mathcal{G}$ to $\mathcal{B}$.
In \eqref{eq:kinVel}, $\Bvdot$ is the rate of change of the components of $\Bv$ in $\mathcal{B}$; the term $\Bw \times \Bv$ accounts for the rotation of the body frame.
These equations hold in 3D for any rigid body, with or without slip.

\autoref{fig:stateGraph} shows which quantities are needed to move from one state to another.
Each arrow lists all the quantities needed to compute the target state; integrating (up) also needs an initial value of the target state.

\begin{figure}[htbp]
  \centering
  \begin{tikzpicture}[
      state/.style  = {draw, rounded corners, minimum width=1.2cm, minimum height=0.8cm},
      sensor/.style = {draw=gray, fill=gray!10, font=\small, align=center},
      calc/.style   = {-{Stealth}, thick},
      meas/.style   = {-{Stealth}, thick, blue!70!black},
      dep/.style    = {-{Stealth}, dashed, gray},
      lbl/.style    = {font=\footnotesize, midway},
    ]

    % --- States grid (rows: x, xdot, xddot; columns: translation, rotation) ---
    \matrix (S) [matrix of math nodes, nodes={state}, row sep=1.2cm, column sep=3cm] {
      |(p)| \Gp & |(R)| \BR \\
      |(v)| \Bv & |(w)| \Bw \\
      |(a)| \Ba &           \\
    };

    % Row and column labels
    \node[anchor=east] at ([xshift=-3.6cm]p.center) {$\bm{x}$};
    \node[anchor=east] at ([xshift=-3.6cm]v.center) {$\dot{\bm{x}}$};
    \node[anchor=east] at ([xshift=-3.6cm]a.center) {$\ddot{\bm{x}}$};
    \node[above=0.3cm of p, font=\small\bfseries] {Translation};
    \node[above=0.3cm of R, font=\small\bfseries] {Rotation};

    % --- Derive / integrate links (down: derive, up: integrate), labelled with the quantities needed ---
    \newcommand{\linkPair}[4]{%
      \draw[calc] ([xshift=-3pt]#1.south) -- node[lbl, left, align=right] {#3} ([xshift=-3pt]#2.north);
      \draw[calc] ([xshift= 3pt]#2.north) -- node[lbl, right, align=left] {#4} ([xshift= 3pt]#1.south);
    }
    \linkPair{p}{v} {$\Gp, \BR$} {$\Bv, \BR$}  % eq:kinPos
    \linkPair{v}{a} {$\Bv, \Bw$} {$\Ba, \Bw$}  % eq:kinVel
    \linkPair{R}{w} {$\BR$}      {$\Bw$}       % eq:kinRot

    % --- Sensor links ---
    \node[sensor, left=0.6cm of p] (ext)     {ICP, GNSS,\\VO, \dots};
    \node[sensor, left=0.6cm of v] (doppler) {Radar\\Doppler};
    \node[sensor, left=0.6cm of a] (acc)     {IMU\\(acc.)};
    \node[sensor, right=0.6cm of w] (gyro)   {IMU\\(gyro)};
    \node[sensor, right=0.6cm of R] (ahrs)   {AHRS};

    \draw[meas] (ext)     -- (p);
    \draw[meas] (doppler) -- (v);
    \draw[meas] (acc)     -- (a);
    \draw[meas] (gyro)    -- (w);
    \draw[meas] (ahrs)    -- (R);

    % AHRS depends on IMU acceleration and angular velocity
    % (routed right of everything drawn so far, so it never crosses a label)
    \coordinate (depX) at ([xshift=0.3cm]current bounding box.east);
    \draw[dep] (gyro.east) -- (gyro.east -| depX) |- (ahrs.east);
    \draw[dep] (acc.south) -- ++(0,-0.5) -| (depX |- ahrs.east) -- (ahrs.east);

    % --- Add new links here ---
    % \draw[dep] (v) -- node[lbl, above] {$\bm{v} = \bm{\omega} \times \bm{r}$} (w);

  \end{tikzpicture}
  \caption{States of the robot. Down arrows differentiate and up arrows integrate following \eqref{eq:kinPos}--\eqref{eq:kinRot}, and their labels give the quantities needed. Blue arrows are direct measurements and dashed arrows are dependencies.}
  \label{fig:stateGraph}
\end{figure}

\subsubsection{Discretization}

Integrating \eqref{eq:kinPos}--\eqref{eq:kinRot} with explicit Euler over a time step $\Delta t$ gives the updates on the left.
Solving them for the lower-order state gives the forward differences on the right, so each pair is an exact inverse.
\begin{align*}
  \Gp[t+\Delta t] & = \Gp[t] + \BRT[t] \, \Bv[t] \, \Delta t
  & \Bv[t] & = \BR[t] \, \frac{\Gp[t+\Delta t] - \Gp[t]}{\Delta t} \\
  \Bv[t+\Delta t] & = \Exp\left(-\Bw[t] \, \Delta t\right) \left(\Bv[t] + \Ba[t] \, \Delta t\right)
  & \Ba[t] & = \frac{\Exp\left(\Bw[t] \, \Delta t\right) \Bv[t+\Delta t] - \Bv[t]}{\Delta t} \\
  \BR[t+\Delta t] & = \Exp\left(-\Bw[t] \, \Delta t\right) \BR[t]
  & \Bw[t] & = -\frac{\Log\left(\BR[t+\Delta t] \, \BRT[t]\right)}{\Delta t}
\end{align*}
$\Exp(\cdot)$ maps a rotation vector to $SO(3)$ and $\Log(\cdot)$ is its inverse.
The position and velocity updates are first order in $\Delta t$.
The rotation update is exact when $\Bw$ is constant over $\Delta t$, as with gyroscope samples.
The forward difference at $t$ is the same value as a backward difference at $t+\Delta t$; only the timestamp differs, by half a step from the true derivative at $t+\Delta t/2$.

Curvature
\begin{align*}
  \kappa & = \frac{\Delta \theta}{\Delta s} \\
  \kappa & = \frac{1}{r}
\end{align*}

Accelerometer model
\begin{align*}
  {}^{\mathcal{B}}\mathbf{\tilde{a}} & = {}^{\mathcal{B}}\mathbf{a} + {}^{\mathcal{B}}_{\mathcal{G}}\mathbf{R} {}^{\mathcal{G}}\mathbf{g} + \mathbf{\eta}(t) + \mathbf{b}(t) \\
  {}^{\mathcal{B}}\mathbf{a} & = {}^{\mathcal{B}}\mathbf{a_c} + \dots \\
  \mathbf{a_c} & = \mathbf{\omega} \times (\mathbf{\omega} \times \mathbf{r'}) \\
  F & = ma
\end{align*}
$\mathbf{r'}$ is the position vector of the object relative to the rotating reference frame. \href{https://en.wikipedia.org/wiki/Coriolis_force}{Wikipedia coriolis force}

Links using scalars
\begin{align*}
  a_c = & \frac{v^2}{r} \\
  v = & \omega  r \\
  v = & \frac{\omega}{\kappa}
\end{align*}

\end{document}
